% Endpoint lemma for the certificates of Theorem 2 (tilted CHSH; Theorem 3 in earlier drafts) (independent audit, 21 Oct 2026).
% Stand-alone file; NOT inserted into prl.tex.
% Text corrected 21 Oct 2026 (referee 2, round 3): (iv) now concludes Phi<=1, and Phi=h=1 only given a tight point.
% No computation changed; the audit was not rerun.  Notation as in prl.tex (Setting, Theorem 1, App. C).
\documentclass[11pt]{article}
\usepackage{amsmath,amssymb,amsthm}
\newtheorem{lemma}{Lemma}
\newtheorem{corollary}{Corollary}
\begin{document}

\noindent\textbf{Notation.} For a functional $g=(a_0,a_1,b_0,b_1,e_{00},e_{01},e_{10},e_{11})$ write
$b=(b_0,b_1)$, $e_x=(e_{x0},e_{x1})$, and for $\delta\in\mathbb R^2$ put $B_{0,1}(\delta)=\frac b2\pm\delta$,
\[
u_a=B_0(\delta)+a\,e_0,\qquad v_{a'}=B_1(\delta)+a'\,e_1\qquad(a,a'=\pm1).
\]
Fix an endpoint $\epsilon\in\{+1,-1\}$ and put $f_\epsilon=(1,\epsilon)$. Let $\mathcal D_\epsilon$ be the
eight deterministic points $L=(A_0,A_1,B_0,B_1)$ with $B_0B_1=\epsilon$, and
\[
h_\epsilon(g)=\max_{L\in\mathcal D_\epsilon} g\cdot L .
\]
For $a,a',\eta\in\{\pm1\}$ define the eight \emph{mixed terms}
\[
M_{a a'\eta}(\delta)=a\,a_0+a'a_1+\eta\,(u_a-v_{a'})\cdot f_\epsilon
 = a\,a_0+a'a_1+\eta\bigl(2\,\delta\cdot f_\epsilon+(a\,e_0-a'e_1)\cdot f_\epsilon\bigr).
\]
For $c\in[-1,1]$ and $s_0,s_1=\pm1$ let
$F_{s_0s_1}(c)=1-s_0a_0-s_1a_1-\|u_{s_0}\|_{G_c}-\|v_{s_1}\|_{G_c}$, so that $\Phi_c(\delta)\le1$ iff
$F_{s_0s_1}(c)\ge0$ for all four sign pairs (App.~C of the paper).

\begin{lemma}[Endpoint lemma]\label{lem:endpoint}
Let $\epsilon\in\{\pm1\}$, $g\in\mathbb R^8$ and $\delta\in\mathbb R^2$ be arbitrary. Then:
\begin{enumerate}
\item[(i)] \emph{(Exact formula.)}
\[
\Phi_\epsilon(\delta)=\max\Bigl\{\,h_\epsilon(g),\ \max_{a,a',\eta}M_{aa'\eta}(\delta)\Bigr\},
\qquad h_\epsilon(g)=h_{K_\epsilon}(g).
\]
\item[(ii)] \emph{(Flat set.)} $\Phi_\epsilon(\delta)=h_{K_\epsilon}(g)$ if and only if
\[
\bigl|\,2\,\delta\cdot f_\epsilon+(a\,e_0-a'e_1)\cdot f_\epsilon\bigr|\ \le\ h_\epsilon(g)-a\,a_0-a'a_1\qquad\text{for all }a,a'=\pm1,
\]
i.e.\ iff the single number $\delta\cdot f_\epsilon$ lies in a (possibly empty) closed interval determined by $g$.
In particular, $\Phi_\epsilon(\delta)$ depends on $\delta$ only through $\delta\cdot f_\epsilon$.
\item[(iii)] \emph{(Signs of the four branches.)} For every sign pair,
\[
F_{s_0s_1}(\epsilon)=\min_{\eta,\eta'=\pm1}\Bigl(1-s_0a_0-s_1a_1-\eta\,u_{s_0}\cdot f_\epsilon-\eta'\,v_{s_1}\cdot f_\epsilon\Bigr),
\]
where the terms with $\eta'=\eta$ equal $1-g\cdot L$ for $L=(s_0,s_1,\eta,\epsilon\eta)\in\mathcal D_\epsilon$ and
the terms with $\eta'=-\eta$ equal $1-M_{s_0s_1\eta}(\delta)$.
\item[(iv)] \emph{(Consequences used in the certificate check.)} Assume
\[
\text{(H1)}\quad g\cdot L\le1\ \ \forall L\in\mathcal D_\epsilon,\qquad
\text{(H2)}\quad M_{aa'\eta}(\delta)<1\ \ \forall a,a',\eta .
\]
Then $\Phi_\epsilon(\delta)\le1$, and if $g\cdot L=1$ for some $L\in\mathcal D_\epsilon$, then
$\Phi_\epsilon(\delta)=h_{K_\epsilon}(g)=1$; $F_{s_0s_1}(\epsilon)\ge0$ for all four pairs; and
$F_{s_0s_1}(\epsilon)=0$ iff $(s_0,s_1,\eta,\epsilon\eta)$ is a tight point ($g\cdot L=1$) for some $\eta$.
For such a tight pair and its $\eta$,
\[
\eta\,u_{s_0}\cdot f_\epsilon>0,\qquad \eta\,v_{s_1}\cdot f_\epsilon>0,
\]
so the radicands $\|u_{s_0}\|^2_{G_c}$, $\|v_{s_1}\|^2_{G_c}$ are positive at $c=\epsilon$, hence on a
neighbourhood of $c=\epsilon$ in $[-1,1]$, and $F_{s_0s_1}$ is real-analytic there.
\end{enumerate}
\end{lemma}

\begin{proof}
For $c=\epsilon=\pm1$ one has $\eta^\top G_\epsilon\eta=\eta_0^2+\eta_1^2+2\epsilon\eta_0\eta_1=(\eta_0+\epsilon\eta_1)^2$,
so $\|\eta\|_{G_\epsilon}=|\eta\cdot f_\epsilon|=\max_{\eta'=\pm1}\eta'\,\eta\cdot f_\epsilon$. Hence
\[
\Phi_\epsilon(\delta)=\max_{a,\eta}\bigl[a\,a_0+\eta\,u_a\cdot f_\epsilon\bigr]+\max_{a',\eta'}\bigl[a'a_1+\eta'\,v_{a'}\cdot f_\epsilon\bigr]
=\max_{a,a',\eta,\eta'}\bigl[a\,a_0+a'a_1+\eta\,u_a\cdot f_\epsilon+\eta'\,v_{a'}\cdot f_\epsilon\bigr].
\]
\emph{Terms with $\eta'=\eta$.} Since $u_a+v_{a'}=b+a\,e_0+a'e_1$ does not depend on $\delta$,
\[
a\,a_0+a'a_1+\eta\,(b+a\,e_0+a'e_1)\cdot f_\epsilon
=a_0A_0+a_1A_1+b_0B_0+b_1B_1+\sum_{x,y}e_{xy}A_xB_y=g\cdot L
\]
with $L=(A_0,A_1,B_0,B_1)=(a,a',\eta,\epsilon\eta)$; as $(a,a',\eta)$ ranges over $\{\pm1\}^3$, $L$ ranges over all of $\mathcal D_\epsilon$.
\emph{Terms with $\eta'=-\eta$.} These are exactly $M_{aa'\eta}(\delta)$, using $u_a-v_{a'}=2\delta+a\,e_0-a'e_1$.
This proves the formula in (i). By Theorem~1 of the paper, $K_\epsilon=\mathrm{conv}\,\mathcal D_\epsilon$, so
$h_{K_\epsilon}(g)=h_\epsilon(g)$.

(ii) By (i), $\Phi_\epsilon(\delta)=h_\epsilon(g)$ iff $M_{aa'\eta}(\delta)\le h_\epsilon(g)$ for all $a,a',\eta$;
taking the maximum over $\eta$ gives the displayed condition, which only involves $\delta\cdot f_\epsilon$.

(iii) $F_{s_0s_1}(\epsilon)=1-s_0a_0-s_1a_1-|u_{s_0}\cdot f_\epsilon|-|v_{s_1}\cdot f_\epsilon|$ and
$|x|+|y|=\max_{\eta,\eta'}(\eta x+\eta' y)$; then argue term by term exactly as in (i).

(iv) By (i), (H1) and (H2), every term in the formula for $\Phi_\epsilon$ is $\le1$, so $\Phi_\epsilon(\delta)\le1$.
If some $L\in\mathcal D_\epsilon$ is tight, then $h_\epsilon(g)=1>\max M_{aa'\eta}$ and (i) gives $\Phi_\epsilon(\delta)=h_{K_\epsilon}(g)=1$.
Without a tight point $\Phi_\epsilon$ may exceed $h_{K_\epsilon}(g)$: for $g=0$, $\delta=(\frac14,0)$ one has
$h_\epsilon(g)=0$, (H1) holds, $M_{aa'\eta}=\pm\frac12<1$, but $\Phi_\epsilon(\delta)=\frac12$. By (iii), $F_{s_0s_1}(\epsilon)$ is the minimum of
the two numbers $1-g\cdot L$ ($L=(s_0,s_1,\eta,\epsilon\eta)$, $\eta=\pm1$), which are $\ge0$ by (H1), and the
two numbers $1-M_{s_0s_1\eta}$, which are $>0$ by (H2); so it vanishes iff one of the two points is tight.
If $L=(s_0,s_1,\eta,\epsilon\eta)$ is tight, then $1-s_0a_0-s_1a_1=\eta\,(u_{s_0}+v_{s_1})\cdot f_\epsilon$. Inserting this into
$M_{s_0s_1\eta}<1$ gives $\eta(u_{s_0}-v_{s_1})\cdot f_\epsilon<\eta(u_{s_0}+v_{s_1})\cdot f_\epsilon$, i.e.\
$\eta\,v_{s_1}\cdot f_\epsilon>0$; inserting it into $M_{s_0s_1,-\eta}<1$ gives $\eta\,u_{s_0}\cdot f_\epsilon>0$.
Finally $\|u\|^2_{G_\epsilon}=(u\cdot f_\epsilon)^2$, and $c\mapsto\|u(c)\|^2_{G_c}$ is a polynomial in $c$ when $\delta=\delta(c)$ is polynomial.
\end{proof}

\begin{corollary}[How the lemma enters Theorem~2]
Let $g$ be a hexagon vertex ($V_1,V_3$ or a $\sigma$-image) and $\delta(c,\theta)$ the certificate~(V13cert) of the paper.
If, for every $\theta\in\Theta$, (H1) holds (it does: these are the hexagon inequalities) and (H2) holds at
$\delta=\delta(\pm1,\theta)$ for both endpoints, then $\Phi_{\pm1}(\delta(\pm1,\theta))=h_{K_{\pm1}}(g)=1$, because each of $V_1,V_3,V_1',V_3'$ has a tight
point in each of $\mathcal D_{+1}$ and $\mathcal D_{-1}$ (exact facet saturations, check (a):
$V_1$: $(1,1,1,1)$, $(1,-1,-1,1)$; $V_3$: $(1,-1,1,1)$, $(-1,1,1,-1)$; $V_1'$: $(1,-1,1,1)$, $(1,1,1,-1)$;
$V_3'$: $(1,1,1,1)$, $(-1,-1,-1,1)$, listed as $(B_0B_1=+1,\ B_0B_1=-1)$), every branch satisfies $F_{s_0s_1}(\pm1,\theta)\ge0$, the zero branches
are exactly the sign pairs $(A_0,A_1)$ of the tight points, and those branches are smooth up to the endpoint.
Hence a box $[1-h,1]\times\Theta_b$ (resp.\ $[-1,-1+h]\times\Theta_b$) is certified by any lower bound showing that
$F$ cannot decrease from its endpoint value $F(\pm1,\theta)\ge0$ into the interior, e.g.\ $\partial_cF<0$
(resp.\ $>0$) on the box, or the second-order bound
$F(c)\ge |c\mp1|\,\bigl(\mp\partial_cF(\pm1)+\tfrac12\min(0,\inf\partial_c^2F)\,h\bigr)$.
Hypothesis (H2) is \emph{not} automatic: it must be verified for the specific $\delta(\pm1,\theta)$
(it is the condition that $\delta(\pm1,\theta)\cdot f_{\pm}$ lies in the flat set of $\Phi_{\pm1}$).
\end{corollary}

\noindent\textbf{Remark (what the paper states).} App.~C of the paper says only ``An endpoint lemma shows that
$\Phi_{\pm1}(\delta(\pm1,\theta))$ equals some $g\cdot L\le1$.'' That statement is true only under (H2); without it,
$\Phi_{\pm1}(\delta)$ can be a mixed term $M_{aa'\eta}$, which is not of the form $g\cdot L$. In the audit,
(H1) and (H2) were verified on all of $\Theta$ for $V_1,V_3,V_1',V_3'$ (all $2\times8$ mixed terms per vertex are $<1$; smallest certified margin $5.8\times10^{-3}$, for $V_3$ at $c=-1$).

\end{document}
